\documentclass[11pt,a4paper]{article}

\usepackage[utf8]{inputenc}
\usepackage[T1]{fontenc}
\usepackage{amsmath,amssymb,amsthm}
\usepackage{graphicx}
\usepackage{xcolor}
\usepackage{hyperref}
\usepackage{cleveref}
\usepackage{algorithm}
\usepackage{algorithmic}
\usepackage{booktabs}
\usepackage{geometry}
\geometry{margin=1in}

% Custom commands
\newcommand{\R}{\mathbb{R}}
\newcommand{\E}{\mathbb{E}}
\newcommand{\norm}[1]{\left\|#1\right\|}

% Theorem environments
\theoremstyle{definition}
\newtheorem{definition}{Definition}
\newtheorem{proposition}{Proposition}
\newtheorem{remark}{Remark}

\title{Action Memory for Smooth Control:\\Preventing Bang-Bang Behavior with Differentiable Zero-Mean Constraints}

\author{Implementation Report\\CaNS Opposition Control with MADDPG}

\date{\today}

\begin{document}

\maketitle

\begin{abstract}
This report addresses two interrelated problems in multi-agent reinforcement learning for turbulent flow control: (1) bang-bang oscillations caused by lack of temporal context, and (2) gradient mismatch from non-differentiable zero-mean constraints. We show that action memory (state augmentation with previous actions) combined with differentiable zero-mean correction enables smooth, physically valid control while maintaining policy trainability. The solution is conceptually simple but requires careful implementation to preserve gradient flow while guaranteeing exact mass conservation.
\end{abstract}

\tableofcontents
\newpage

\section{Problem Statement}

\subsection{The Bang-Bang Control Problem}

Multi-agent RL policies for flow control exhibit high-frequency oscillations (bang-bang behavior) characterized by:
\begin{itemize}
    \item Rapid switching between extreme actions: $a_t \in \{-1, +1\}$
    \item High temporal gradient: $\norm{\nabla_t a}_{\text{RMS}} \approx 0.4$
    \item Good drag reduction but poor physical realism
\end{itemize}

\textbf{Root cause}: The policy $\pi(\mathbf{u}_t, \mathbf{w}_t)$ creates a closed-loop feedback instability:
\begin{equation}
    \text{Bang-bang } a_t \to \text{Flow waves} \to \text{Oscillating } (\mathbf{u}_{t+1}, \mathbf{w}_{t+1}) \to \text{Bang-bang } a_{t+1}
\end{equation}

Without temporal context, the policy cannot detect or dampen this self-reinforcing cycle.

\subsection{The Zero-Mean Constraint Problem}

Flow control requires mass conservation: $\sum_{i,j} a_{ij} = 0$. Standard implementations apply this as \textbf{post-processing}:
\begin{equation}
    a^{\text{exec}} = a^{\text{raw}} - \text{mean}(a^{\text{raw}})
\end{equation}

\textbf{Critical issues}:
\begin{enumerate}
    \item \textbf{Gradient mismatch}: Backpropagation doesn't see the correction, policy assumes direct control
    \item \textbf{Induced oscillations}: Varying mean $\mu_t$ creates oscillating corrections even for smooth $a^{\text{raw}}_t$
    \item \textbf{Credit assignment failure}: Policy never learns to output near-zero-mean actions
\end{enumerate}

\section{Unified Solution: Action Memory + Differentiable Constraints}

\subsection{State Augmentation with Previous Action}

\textbf{Modified state representation}:
\begin{equation}
    \mathbf{s}_t = [\mathbf{u}_t, \mathbf{w}_t, \mathbf{a}_{t-1}] \quad \in \R^{3 \times 8 \times 8}
\end{equation}

\textbf{Benefits}:
\begin{itemize}
    \item Policy sees full feedback loop: $\pi(\mathbf{u}_t, \mathbf{w}_t, \mathbf{a}_{t-1})$ learns "I was oscillating, dampen response"
    \item Enables temporal consistency without explicit penalties
    \item Maintains Markov property (history summarized in $\mathbf{a}_{t-1}$)
\end{itemize}

\subsection{Differentiable Zero-Mean Correction}

\textbf{Key insight}: Make the zero-mean correction part of the computational graph during training.

\subsubsection{Training Forward Pass}
\begin{algorithm}
\caption{Differentiable Zero-Mean Correction (Training)}
\begin{algorithmic}[1]
\STATE \textbf{Forward pass}: $\mathbf{A}^{\text{raw}} = \pi_\theta(\mathbf{S})$ \quad [$\in [-1,1]$ from $\tanh$]
\STATE \textbf{Zero-mean correction (differentiable)}:
\[
\mu = \text{mean}(\mathbf{A}^{\text{raw}}), \quad \mathbf{A}^{\text{zm}} = \mathbf{A}^{\text{raw}} - \mu
\]
\STATE \textbf{Add exploration noise}: $\boldsymbol{\epsilon} \sim \mathcal{N}(0, \sigma^2 I)$, $\boldsymbol{\epsilon} \leftarrow \boldsymbol{\epsilon} - \text{mean}(\boldsymbol{\epsilon})$
\STATE \textbf{Add noise}: $\mathbf{A}^{\text{noisy}} = \mathbf{A}^{\text{zm}} + \boldsymbol{\epsilon}$
\STATE \textbf{Clip to bounds}: $\mathbf{A}^{\text{clip}} = \text{clip}(\mathbf{A}^{\text{noisy}}, -1, 1)$ \quad [breaks zero-mean!]
\STATE \textbf{Final zero-mean guarantee}: $\mathbf{A} = \mathbf{A}^{\text{clip}} - \text{mean}(\mathbf{A}^{\text{clip}})$
\STATE \textbf{Compute loss}: $\mathcal{L} = -Q(\mathbf{S}, \mathbf{A}) + \lambda_s \mathcal{L}_{\text{spatial}} + \lambda_g \mathcal{L}_{\text{global\_mean}}$
\STATE \textbf{Backpropagate}: $\nabla_\theta \mathcal{L}$ \quad [gradients flow through all corrections]
\end{algorithmic}
\end{algorithm}

\subsubsection{Inference (No Noise)}
\begin{algorithm}
\caption{Inference with Exact Zero-Mean}
\begin{algorithmic}[1]
\STATE \textbf{Forward pass}: $\mathbf{A}^{\text{raw}} = \pi_\theta(\mathbf{S})$ \quad [$\in [-1,1]$ from $\tanh$]
\STATE \textbf{Zero-mean guarantee}: $\mathbf{A} = \mathbf{A}^{\text{raw}} - \text{mean}(\mathbf{A}^{\text{raw}})$
\STATE \textbf{Execute}: $\mathbf{A}$ is exactly zero-mean, no clipping needed
\end{algorithmic}
\end{algorithm}

\section{How Trainability and Enforcement Coexist}

\subsection{The Seeming Paradox}

\textbf{Question}: If we always apply $a = a^{\text{raw}} - \text{mean}(a^{\text{raw}})$, how does the policy learn anything? Won't it just output arbitrary actions since they'll be corrected anyway?

\textbf{Answer}: Differentibility makes all the difference.

\subsection{Gradient Flow Explanation}

Consider a single training step with 2 actions for simplicity: $a^{\text{raw}} = [a_1, a_2]$.

\subsubsection{Without Differentiable Correction (Post-Processing)}
\begin{align}
    \text{Forward:} \quad &a^{\text{raw}} = [0.5, -0.3] \quad \Rightarrow \quad \mu = 0.1 \\
    &a^{\text{exec}} = [0.5 - 0.1, -0.3 - 0.1] = [0.4, -0.4] \\
    \text{Gradient:} \quad &\frac{\partial \mathcal{L}}{\partial a_1} = -2.0 \quad (\text{from critic feedback}) \\
    \text{Update:} \quad &\theta \leftarrow \theta + \eta \cdot 2.0 \quad (\text{policy doesn't see correction!})
\end{align}

Policy learns as if it directly controls $[0.5, -0.3]$, but environment saw $[0.4, -0.4]$. \textbf{Mismatch!}

\subsubsection{With Differentiable Correction}
\begin{align}
    \text{Forward:} \quad &a^{\text{raw}} = [0.5, -0.3], \quad \mu = \text{mean}([0.5, -0.3]) = 0.1 \\
    &a = a^{\text{raw}} - \mu = [0.4, -0.4] \\
    \text{Gradient chain rule:} \quad &\frac{\partial \mathcal{L}}{\partial \theta} = \frac{\partial \mathcal{L}}{\partial a} \cdot \frac{\partial a}{\partial a^{\text{raw}}} \cdot \frac{\partial a^{\text{raw}}}{\partial \theta}
\end{align}

The Jacobian captures coupling:
\begin{equation}
    \frac{\partial a_i}{\partial a_j^{\text{raw}}} = \begin{cases}
        1 - \frac{1}{N} & i = j \\
        -\frac{1}{N} & i \neq j
    \end{cases}
\end{equation}

\textbf{Key effect}: If $a_1$ increases, $a_2$ decreases (to maintain zero-mean). Policy learns this coupling!

\subsection{Learning Dynamics}

\textbf{Early training}:
\begin{itemize}
    \item Policy outputs arbitrary $a^{\text{raw}}$ with large $|\mu|$ (e.g., $\mu = 0.5$)
    \item Large corrections applied: $a = a^{\text{raw}} - 0.5$
    \item Gradients penalize large mean through:
    \begin{equation}
        \mathcal{L}_{\text{global\_mean}} = \lambda_g \left(\text{mean}(a^{\text{raw}})\right)^2
    \end{equation}
\end{itemize}

\textbf{Late training}:
\begin{itemize}
    \item Policy learns to output $a^{\text{raw}} \approx \text{zero-mean}$
    \item Small corrections: $|\mu| < 0.05$
    \item Actions are \textbf{inherently near-zero-mean}, not arbitrarily corrected
    \item Final correction is just a safety guarantee
\end{itemize}

\subsection{Why This Doesn't Break Learning}

\textbf{Intuition}: The constraint $\sum a_i = 0$ removes \textbf{one degree of freedom} (the global offset). But we have $N = 64 \times 64 = 4096$ actions. The policy still has 4095 degrees of freedom to optimize!

\textbf{Analogy}: Imagine training a robot arm to reach targets while constrained to a table surface. The vertical position is always $z = 0$ (enforced). Does this prevent learning? No! The robot learns to position $(x, y)$ optimally. Similarly, the policy learns the \textit{pattern} of actions (spatial structure) while the global mean is constrained to zero.

\section{Critical Implementation Details}

\subsection{Noise and Clipping Interaction}

\textbf{Problem}: Clipping breaks zero-mean!
\begin{align}
    a^{\text{noisy}} &= [1.2, -0.8, 0.1, -0.5] \quad (\text{mean} = 0) \\
    a^{\text{clip}} &= [1.0, -0.8, 0.1, -0.5] \quad (\text{mean} = -0.05 \neq 0)
\end{align}

\textbf{Solution}: Always re-apply zero-mean after clipping:
\begin{equation}
    a^{\text{final}} = a^{\text{clip}} - \text{mean}(a^{\text{clip}})
\end{equation}

\textbf{Optimization}: When noise is annealed to zero ($\sigma = 0$ in late training), skip clipping entirely since $\tanh$ already guarantees $a \in [-1, 1]$.

\subsection{Global vs Per-Tile Zero-Mean}

\textbf{Invalid approach}: Forcing each $8 \times 8$ agent tile to have zero-mean:
\begin{equation}
    \mathcal{L}_{\text{per-tile}} = \sum_{i=1}^{64} \left(\text{mean}(\mathbf{a}_i)\right)^2 \quad \textcolor{red}{\text{WRONG!}}
\end{equation}

\textbf{Why invalid}: Different flow regions require different average actuation:
\begin{itemize}
    \item Upstream: net blowing $(\text{mean} > 0)$
    \item Downstream: net suction $(\text{mean} < 0)$
\end{itemize}

\textbf{Correct approach}: Global zero-mean only:
\begin{equation}
    \mathcal{L}_{\text{global}} = \left(\frac{1}{64 \times 64}\sum_{i,j} a_{ij}\right)^2
\end{equation}

This respects physics (mass conservation) without artificial constraints.

\section{Results and Validation}

\subsection{Expected Training Behavior}

\textbf{Episode 1-3 (Buffer warmup)}:
\begin{itemize}
    \item Random exploration, filling replay buffer
    \item No training yet ($\texttt{training\_step} = 0$)
    \item Only basic metrics logged
\end{itemize}

\textbf{Episode 4+ (Training active)}:
\begin{itemize}
    \item $\texttt{global\_mean\_loss}$ decreases: Policy learning constraint
    \item $\texttt{action\_mean}$ approaches zero: $|\mu| < 0.05$
    \item $\texttt{actor\_grad\_norm}$ stable: $< 10$
    \item $\texttt{dpdx\_mean}$ decreases: Drag reduction improving
\end{itemize}

\subsection{Convergence Indicators}

\begin{center}
\begin{tabular}{lll}
\toprule
\textbf{Metric} & \textbf{Convergence Target} & \textbf{Meaning} \\
\midrule
$\texttt{global\_mean\_loss}$ & $< 0.01$ & Policy outputs near-zero mean \\
$\texttt{action\_mean}$ & $< 0.05$ & Small corrections needed \\
$\texttt{actor\_exploding}$ & $= 0$ & Stable training \\
$\texttt{temporal\_gradient\_rms}$ & $< 0.2$ & Smooth actions \\
\bottomrule
\end{tabular}
\end{center}

\subsection{Action Memory Impact}

\textbf{Without action memory}:
\begin{itemize}
    \item Temporal gradient RMS: $\approx 0.4$
    \item FFT shows high-frequency peaks
    \item Bang-bang oscillations persist
\end{itemize}

\textbf{With action memory}:
\begin{itemize}
    \item Temporal gradient RMS: $< 0.2$
    \item Smooth frequency spectrum
    \item Adaptive temporal consistency
\end{itemize}

\section{Summary: Key Insights}

\begin{enumerate}
    \item \textbf{Action memory breaks feedback loops}: Seeing $\mathbf{a}_{t-1}$ lets policy detect and dampen oscillations

    \item \textbf{Differentiability $\neq$ no constraint}: Making the correction differentiable allows gradient flow while still enforcing exact zero-mean

    \item \textbf{Policy learns the constraint}: Through gradient coupling, network learns to output near-zero-mean actions, minimizing correction magnitude

    \item \textbf{Safety + trainability}: Differentiable correction during training + final guarantee ensures both learning and physical validity

    \item \textbf{Global constraint only}: Per-tile zero-mean is physically invalid; only global mass conservation matters

    \item \textbf{Clipping requires re-correction}: Always reapply zero-mean after clipping to guarantee exact constraint satisfaction
\end{enumerate}

\section{Code Implementation}

\subsection{Training Script Pattern}

\begin{verbatim}
# Forward pass through all agents
actions_raw = maddpg.act(observations)  # tanh \in [-1, 1]

# Differentiable zero-mean correction
global_mean = actions_raw.mean()
actions = actions_raw - global_mean  # Gradients flow!

# Add zero-mean exploration noise (if training)
if noise_scale > 0:
    noise = np.random.normal(0, noise_scale, actions.shape)
    noise = noise - noise.mean()  # Zero-mean noise
    actions = actions + noise
    actions = np.clip(actions, -1, 1)  # May break zero-mean

    # CRITICAL: Re-apply zero-mean after clipping
    actions = actions - actions.mean()

# actions is now exactly zero-mean, ready for environment
next_observations, rewards, dones, infos = env.step(actions)

# Update networks with gradient flow through zero-mean correction
maddpg.update_batched(batch)
\end{verbatim}

\subsection{Loss Components}

\begin{align}
    \mathcal{L}_{\text{actor}} &= -\E[Q(\mathbf{s}, \mathbf{a})] + \lambda_s \mathcal{L}_{\text{spatial}} + \lambda_g \mathcal{L}_{\text{global\_mean}} + \lambda_c \mathcal{L}_{\text{consistency}} \\
    \mathcal{L}_{\text{spatial}} &= \E\left[\norm{\nabla_x \mathbf{a}}^2 + \norm{\nabla_y \mathbf{a}}^2\right] \\
    \mathcal{L}_{\text{global\_mean}} &= \left(\text{mean}(\mathbf{a})\right)^2 \\
    \mathcal{L}_{\text{consistency}} &= \E\left[\sum_{\text{similar pairs}} \max(0, \norm{\mathbf{a}_i - \mathbf{a}_j} - m)\right]
\end{align}

\textbf{Recommended hyperparameters}:
\begin{itemize}
    \item $\lambda_s = 0.5$ (spatial smoothness)
    \item $\lambda_g = 0.05$ (global zero-mean)
    \item $\lambda_c = 0.0$ (consistency - often not needed with action memory)
\end{itemize}

\section{Conclusion}

The combination of action memory and differentiable zero-mean correction solves two problems simultaneously:
\begin{enumerate}
    \item \textbf{Bang-bang oscillations}: Action memory provides temporal context to detect and prevent feedback instability
    \item \textbf{Constraint learning}: Differentiable correction enables policy to learn the zero-mean structure rather than fighting against it
\end{enumerate}

The key conceptual insight is that \textbf{differentiability enables learning of the constraint}, not removal of the constraint. The policy learns to operate within the zero-mean manifold, making corrections progressively smaller while still maintaining exact mass conservation for solver stability.

\end{document}
